\setcounter{section}{29}
\section{Gobernanza Tic}

Nombre completo: Modelos de gobernanza TIC. Organización e instrumentos operativos de las tecnologías de la información y las comunicaciones en la Administración General del Estado y sus organismos públicos. La transformación digital de la Administración General del Estado.

\noindent\fbox{\parbox{\textwidth}{Apuntes actualizados al año 2019}}

\localtableofcontents

\subsection{Introducción}
Junto con las leyes 39/2015 y 40/2015 se configura una administración pública y una interacción con la ciudadanía eminentemente digital. Es inevitable el protagonismo de las TIC en la vida, y la administración pública no es menos. La transformación digital juega un papel importante para las organizaciones a la hora de optimizar su funcionamiento. Con la cantidad de organizaciones que existen enfrentándose a este reto simultaneamente, es seguro decir que es una actividad importante.

\subsection{Modelos de gobernanza TIC y Su implantación en las administraciones públicas}
    \subsubsection{La gobernanza en la AGE}
        Tradicionalmente se ha desarrollado en las unidades de informática de los distintos ministerios. Con autonomía para la gestión de los procesos, tanto en adquisiciones, como en el desarrollo y mantenimiento. De esta manera se obtiene independencia y agilidad, y mejor conociemiento del negocio. Por contra, esta descentralización incrementa la dificultad de coordinación interdepartamental,  la posibilidad de la economía de escala y dificulta la reutilización de servicios o productos desarrollados en los departamentos.

        Para paliar la situación se han creado órganos colegiados como el Consejo Superior de Informática, que pasó a llamarse Consejo Superior de Administración Electrónica, siendo éste el órgano máximo de en materia de administración electrónica, de donde han emanado las líneas y proyectos de administración electrónica.

        También se han creado diferentes órganos, que han permitido realizar el seguimiento y control de las diferentes inversiones y gastos a nivel ministerial. Esto ha impedido que, salvo en algunos ministerios, el desarrollo de una estrategia digital que soporte los procesos administrativos.

        En 2012, la CORA (Comisión para la Reforma de las Administraciones públicas) hace un estudio de los modelos de gobernanza en las TIC. Teniendo en cuenta lo realizado por algunas comunidades autónomas, otros países europeos y empresas internacionales. Se decide que el modelo ha de ser el siguiente:

    
        \noindent\fbox{\parbox{\textwidth}{
            \textit{(…) para llevar a cabo la estrategia propuesta es necesario proceder a un rediseño de la organización actual creando una estructura de gobernanza que facilite la obtención de los objetivos de consolidación y la gestión de los medios comunes. En este modelo se propone la creación en un primer momento de la figura del CIO de la AGE como máximo responsable de la Administración en el ámbito de las TIC y con posterioridad ha de valorarse la creación de una Agencia como organismo encargado de consolidar y desarrollar un catálogo de servicios horizontales para la prestación de los servicios comunes a toda la Administración, así como de racionalizar los desarrollos de infraestructuras y servicios sectoriales. Mediante este modelo se permitirá lograr economías de escala por vía de la consolidación de infraestructuras, servicios y compras.
            Se propone, por tanto, la creación de un órgano TIC para toda la Administración General del Estado al máximo nivel (CIOAGE), que pueda gestionar el desarrollo de un catálogo de servicios horizontales para la prestación de los servicios comunes, además de racionalizar las infraestructuras y servicios sectoriales.}
        }}

        Con esta propuesta entra en vigor el RD 695/2013 por el que se crea la Dirección de Tecnologías de la Información y las Comunicaciones (DTIC) de la Administración General del Estado. Con Rango de subsecretaría y dependiendo de los ministerios de Presidencia y Hacienda y Administraciones Públicas. al frente de esta se encontraba el Director de Tecnologías de Información y las Comunicaciones (CIOAGE). Este rango se le da para dotarles de autoridad para llevar a cabo los cambios para el nuevo modelo organizativo.
        
        Actualmente esta figura la desempeña el Secretario General de la Administración Digital, según el Art. 7 del RD 863/2018 Corresponde a la SGAD el ejercicio de las siguientes funciones:
        \begin{itemize}
            \item Elaboración de estrategia de Administración Digital y Servicios Públicos Digitales de la AGE
            \item Gestión compartida de los servicios comunes de Sistemas de Información y Comunicación para la AGE
            \item Diseño técnico, implantación y gestión de los medios y servicios necesarios para evolucionar los servicios públicos actuales
            \item Racionalización e impulso de integración tecnológica de los servicios de atención ciudadana
            \item Supervisar la ejecución de medidas específicas en los planes de acción departamentales para la transformación digital
            \item Impulsar la consolidación de servicios, infraestructuras, equipamientos y redes informáticas comunes de la AGE
            \item Realización de propuestas e interlocución con organismos especializados en materias de seguridad y protección de datos. Como el Centro Criptológico Nacional
            \item Ser referente nacional en el ámbito de competencias digitales y presidir la Comisión Sectorial de Administración Electrónica
            \item Definición de estándares y directrices técnicas y de gobierno TIC, normas de seguridad y calidad
            \item Elaboración, desarrollo, implantación y gestion del Catálogo de Servicios Compartidos
            \item \textit{et al}
        \end{itemize}
        
        El RD 806/2014 describe el modelo de gobernanza TIC con tres principios básicos:
        \begin{description}
            \item[Disponer de una planificación estratégica común] Orientando las actuaciones y líneas estratégicas de las TIC
            \item[Potenciar la Administración digital y las TIC] como instrumentos para hacer sostenible el proceso de innovación y la calidad de los servicios
            \item[Racionalizar el uso de los recursos informáticos] en busca de la eficiencia, ahorrando costes, especialmente en la actividad administrativa, especialmente por la homogeneidad y simplicidad
        \end{description}
        Para ello se reestructuraron los óírganos colegiados, extinguiéndose el CSAE y CMAE, y constituyéndose nuevos órganos. En lugar de tener que coordinarse estrategias, se centraliza una única de caracter horizontal.

    \subsubsection{La Gobernanza a nivel nacional}
        España es un país complejo con 3 niveles administrativos independientes, amén de otras entidades locales. El modelo de gobernanza debe permitir alinear las acciones de los actores involucrados. Trabajando por consensos.

        La colaboración TIC se circunscribe a un órgano de competencias ámplias: la Conferencia Sectorial de la Administración Pública (CSAP). Encargado de la cooperación entre la AGE, Administraciones autónomas y locales. Funcionando según el Reglamento de Organización y Funcionamiento. Esta regulación parte de lo especificado en la ley 40/2015. Siendo las funciones de la comisión sectorial las siguientes:
        \begin{itemize}
            \item Asegurar la compatibilidad e interoperabilidad de los sistemas y aplicaciones
            \item Impulsar el desarrollo de la administración electrónica
            \item Asegurar la cooperación entre administraciones públicas
        \end{itemize}

        Se establecen grupos de trabajo técnicos para compartir experiencias, buenas prácticas y realizar proyectos de colaboración. Algunos grupos son:
        \begin{itemize}
            \item Identidad Digital y Firma Electrónica
            \item de Intermediación y Servicios. Carpeta ciudadana
            \item Interoperabilidad
            \item Telecomunicaciones integradas
            \item Observatorio, Indicadores y Medidas
            \item Seguridad
            \item Reutilización de Aplicaciones y Servicios y RISP
            \item Repositorios Comunes
            \item Documento, Expediente y Archivo Electrónico
            \item Funcionarios Habilitados y Registro Electrónico de Apoderamientos
            \item Apoyo funcional a la implantación de las leyes 39 y 40/2015
            \item Identificador Europeo de Legislación
            \item Factura Electrónica
        \end{itemize}

\subsection{RD 806/2014 Sobre organización e instrumentos operativos de las TIC en la AGE}
    Ciertas referencias del RD han sido modificadas posteriormente, algunas sustancialmmente.

    \subsubsection{Comisión de Estrategia TIC - CETIC}
        \subparagraph{Funciones de la Comisión de Estrategia TIC}
            Sus funciones están descritas en el art. 4 del RD. Y son:
            \begin{itemize}
                \item Fijar líneas estratégicas en materia de tecnologías de la información y comunicaciones
                \item Aprobar la propuesta de Estrategia TIC de la AGE para su elevación a consejo de ministros.
                \item Informar anteproyectos de ley, de disposiciones reglamentarias y otras normas propuestas para regular en materias TIC
                \item Definir prioridades de inversión en materias TIC
                \item Declarar los medios o servicios compartidos
                \item Declarar los proyectos de interés prioritarios a propuesta de ministerios y organismos públicos
                \item Impulsar la colaboración y cooperación con comunidades autónomas y entidades locales
                \item Impulsar la cooperación de la AGE con la UE, organizaciones internacionales y con Iberoamérica
                \item Actuar como observatorio de la Administración Electrónica y Transformación Digital
            \end{itemize}

        \subparagraph{Organización de la comisión de Estrategia TIC}
            La Comisión tiene un pleno integrado por los titulares de las secretarías de Estado de función pública, para el Avance Digital y de la Seguridad Social, y por los subsecretarios, o titulares de un órgano superior de los departamentos ministeriales. Es presidido por el Ministro de Política Territorial y Función Pública. Se deberían reunir dos veces al año.

    \subsection{Comisiones Ministeriales de la Administración Digital}
        Son órganos colegiados de ámbito departamental, respomsables del impulso y coordinación interna en cada departamento en materia de Administración Digital, y son óírganos de enlace con la SGAD.

        Las CMAD  estudian y planifican las necesidades funcionales de las áreas administrativas del ministerio. Desempeñan las funciones de:
        \begin{itemize}
            \item Actuar como órgano de relación entre departamentos ministeriales para asegurar la coordinación con criterios y políticas definidas por la AGE
            \item Impulsa, ejecuta y supervisa en el ámbito del departamento, el cumplimiento de directrices y seguimiento de pautas de actuación, según la Estrategia TIC de la AGE
            \item Elabora el plan de acción del departamento para la transformación difital según los criterios establecidos por la SGED
            \item Analiza necesidades funcionales de las unidades de gestión del departamento y organismos adscritos, evaluando alternativas de solución propuestas, identificando oportunidades de mejora y estimando costes en recursos humanos y materiales
            \item Impulsa la digitalización de servicios y procedimientos del departamento para homogeneizarlos
            \item Colaborar con la SGAD en la identificación y puesta a disposición de los medios humanos, materiales y económicos para la puesta en marcha y funcionamiento de los servicios compartidos.
        \end{itemize}
    \subsection{Comité de Dirección de las Tecnologías de la Información y las Comunidades}
        Órgano adscrito a la SGAD, que se reune mensualmente. Integrado por el responsable TIC de las subsecretarías del órgano superior correspondiente. Actúa como óírgano de coordinación y colaboración entre la Secretaría General de Administración Digital y los órganos integrantes de la AGE y organismos públicos a fin de establecer una acción coordinada de acuerdo con las líneas definidas por la Comisión de estrategia TIC.

    \subsection{Unidades TIC}
        Son aquellas que, fuera del ámbito de los órganos colegiados, su función sea la provisión de servicios en materias TIC. Sus funciones deben ser una o varias de las siguientes:
        \begin{itemize}
            \item Soporte, operación, implementación o gestión de sistemas informáticos
            \item Desarrollo de aplicativos informáticos en entornos multiusuario
            \item Consultoría informática
            \item Seguridad
            \item Atención técnica a usuarios
            \item Innovación en el ámbito de las tic
            \item Administración digital
            \item Confirmar la voluntad de adquisición de bienes y servicios en el ámbito de las TIC
            \item Otras funciones relevantes en materia TIC
        \end{itemize}

\subsection{La transformación digital de la administraicón general del estado}
    El Plan de Transformación digital de la Administración General del Estado y sus Organismos Públicos constituye el marco estratégico para avanzar en la transformación de la Administración. Establece principios rectores, objetivos y acciones para alcanzarlo. También establece ciertos hitos.

    Se incorporan las recomendaciones de la OCDE para elaborar estrategias de administración digital. Se alinean con los servicios, actos y políticas de la Unión Europea.  y la nueva estrategia de la Comisión Europea para el Mercado Único Digital. También es un instrumento para aplicar las leyes 39 y 40 de 2015.

    Se impulsó por la DTIC, pero ha sido un esfuerzo colaboratrivo donde han participado todos los departamentos de la AGE.

    Los principios rectores son
    \begin{itemize}
        \item Vocación de orientación al usuario del servicio
        \item Unidad y visión integral de las necesidades y acciones
        \item Asegurar la colaboración inter e intra administrativa y el fomento de Alianzas. 
        \item Transparencia y rendimiento de cuentas
        \item Innovación para impulsar el progreso y adaptación a los cambios.
    \end{itemize}

    También se establecen objetivos estratégicos con ciertas líneas de acción asociadas.
    \begin{description}
        \item[Incremento de la productividad y eficacia en el funcionamiento interno de la Administración] Transformar los procesos internos en electrónicos y desarrollar el puesto de trabajo digital
        \item[Canal digital como medio preferido para relacionarse con la Administración] Proveer servicios públicos digitales adaptaos a las nuevas tecnologíás. Mejorar satisfacción del usuario  de servicios públicos digitales. Proveer la innovación en la prestación de servicios
        \item[Mayor eficiencia en la prestación de servicios TIC] Proveer de manera compartida servicios comunes
        \item[Gestión corporativa inteligente del conocimiento, informacion y datos] Publicar información disponible. Disponer de sistemas de análisis de datos para la toma de decisiones
        \item[Estrategia corporativa de seguridad y usabilidad] Garantizar la seguridad de los sistemas de información.
    \end{description}
    Algunos servicios compartidos son:
    \begin{itemize}
        \item Rube híbrida (red SARA)
        \item Correo electrónico unificado
        \item servicio multicanal de atención al ciudadano
        \item gestión de notificaciones
        \item gestión de expediente y documento electrónico
        \item gestión de archivo electrónico
    \end{itemize}